\documentclass[11pt]{report}

\usepackage{geometry}
\geometry{letterpaper, portrait, left=1.5in, right=1in, top=1.5in, bottom=1in}

\usepackage{lmodern}
\usepackage{lscape}

\usepackage{etoolbox}
\usepackage[style=ieee, citestyle=numeric-comp, sortcites, backend=biber, defernumbers=true, url=false]{biblatex}
\addbibresource{references.bib}

% Prevent page breaks within bibliography entries
\patchcmd{\bibsetup}{\interlinepenalty=5000}{\interlinepenalty=10000}{}{}
\setcounter{biburlnumpenalty}{100}
\setcounter{biburlucpenalty}{100}
\setcounter{biburllcpenalty}{100}

\setlength{\parindent}{0pt}
\setlength{\parskip}{0.5\baselineskip}
\usepackage{titlesec}
% Set consistent spacing before and after section headers
\titleformat{\chapter}[display]
  {\normalfont\large\bfseries\centering}{\chaptertitlename\ \thechapter}{-0.5em}{\large\centering}
\titleformat{\section}
  {\normalfont\normalsize\bfseries}{\thesection}{1em}{}
\titleformat{\subsection}
  {\normalfont\normalsize\bfseries}{\thesubsection}{1em}{}
\titleformat{\subsubsection}
  {\normalfont\normalsize\bfseries}{\thesubsubsection}{1em}{}
\titlespacing*{\chapter}{0pt}{0.25\baselineskip}{0.25\baselineskip}
\titlespacing*{\section}{0pt}{0.25\baselineskip}{0.25\baselineskip}
\titlespacing*{\subsection}{0pt}{0.25\baselineskip}{0.25\baselineskip}
\titlespacing*{\subsubsection}{0pt}{0.25\baselineskip}{0.25\baselineskip}

\usepackage[separate-uncertainty=true, multi-part-units=single]{siunitx}
\sisetup{
	exponent-to-prefix = true,
	range-independent-prefix = true, 
  range-units = single,
  % range-phrase = {-},
  group-separator = {,},
  group-digits = integer,
  display-per-mode = fraction, 
  inline-per-mode = symbol,
  group-minimum-digits = 4
}

% \DeclareSIUnit\second{second}
% \DeclareSIUnit\minute{minute}
% \DeclareSIUnit\hour{hour}
% \DeclareSIUnit\day{day}

\usepackage{minted}
\usemintedstyle{friendly}
\setminted{
  frame=single,
  framesep=2mm,
  baselinestretch=1.0,
  breaklines=true,
  breakanywhere=true,
  fontsize=\footnotesize,
}

\usepackage{tabularx}
\usepackage{booktabs}
\usepackage{multirow}
\usepackage{graphicx}
\usepackage{xurl}
% \usepackage{subfig} % REMOVE this line to avoid conflict with subcaption
\usepackage{subcaption}
\usepackage{caption}
\captionsetup{font=footnotesize, margin=10pt, skip=5pt} % Set caption font to 8pt with smaller margins
\usepackage{amsmath} % Added for equation* environment
\usepackage{hyperref}
% Fast compile mode: set to \fastcompilefalse when actively reviewing TODOs/layout.
\newif\iffastcompile
\fastcompiletrue

% \iffastcompile
% \usepackage[disable]{todonotes}
% \else
\usepackage{todonotes}
% \usepackage{showframe} % Shows page layout boxes - keep this off unless debugging layout
% \fi
\usepackage{setspace}
\usepackage{calc}

\usepackage{fancyhdr}
\pagestyle{fancy}
\fancyhf{}
% this prints a line under the header
\renewcommand{\headrulewidth}{0 pt}
%this prints a line under the footer
\renewcommand{\footrulewidth}{0 pt}
\fancyhead[RO]{}
\fancyhead[LO]{}
\fancyfoot[C]{}
\rhead{\thepage}

\fancypagestyle{plain}{%
	\fancyhf{}
	\rhead{\thepage}
}

\usepackage{pdfpages}    % For including PDF files
\usepackage[toc,page]{appendix}  % For structured appendices
\usepackage{subfiles}  % For subfile support


\title{VIBRO-ACOUSTIC SIGNATURES AND DETECTION ALGORITHMS FOR STORED PRODUCT INSECTS}
\author{Daniel Kadyrov, Ph.D. Candidate}
\date{\today}

\newcommand{\thesistitle}{VIBRO-ACOUSTIC SIGNATURES AND DETECTION ALGORITHMS FOR STORED PRODUCT INSECTS}
\newcommand{\thesisauthor}{Daniel Kadyrov}
\newcommand{\thesisadvisor}{Raju Datla}
\newcommand{\thesisyear}{2025}
\newcommand{\thesisname}{Daniel Kadyrov}
\newcommand{\thesischairadvisor}{Dr. Raju Datla}    % this name prints on the title page as chairman and the abstract page as advisor
\newcommand{\committeenameA}{Dr. Alexander Sutin, Co-Chairperson}
\newcommand{\committeenameB}{Dr. Muhammad R. Hajj}
\newcommand{\committeenameC}{Dr. Richard Mankin}
\newcommand{\thesisdepartment}{Department of Civil, Environmental and Ocean Engineering}
\newcommand{\thesisdate}{November, 2025}
\newcommand{\thesistype}{dissertation}
\newcommand{\thesisdegree}{Doctor of Philosophy}
\newcommand{\graddate}{\the\year} % like 2020, 2019, no month or day 

\makeatletter
\renewcommand{\maketitle}{
\begin{center}
  {\thesistitle}
  \vspace{.15in}
  
    by
    
  \vspace{.15in}
  \thesisauthor
  
  \vspace{.15in}
  
 {A DISSERTATION}\\
  \vspace{.2in}
  \begin{spacing}{1}
    {Submitted to the Faculty of the Stevens Institute of Technology\\
    in partial fulfillment of the requirements for the degree of}
    \end{spacing}
  \vspace{.2in}
  
  {DOCTOR OF PHILOSOPHY}\\
  \vspace{1.0in}
  % for master thesis, change Chairman to Advisor
    \hfill 
    \begin{minipage}{80mm}
    \begin{spacing}{ }\noindent \rule{3.2in}{0.1mm}
        \thesisname, Candidate\\[3mm]
        \underline{ADVISORY COMMITTEE}\\[3mm]
        \noindent \rule{3.2in}{0.1mm}\\[-1.3mm]
        % for master thesis, change Chairman to Advisor
        {\thesischairadvisor}, Chairperson  \hfill{Date}\\[2mm]
        {\noindent \rule{3.2in}{0.1mm}}\\[-1.3mm]
        {\committeenameA}        \hfill{Date}\\[2mm]
        {\noindent \rule{3.2in}{0.1mm}}\\[-1.3mm]
        {\committeenameB}        \hfill{Date}\\[2mm]
        {\noindent \rule{3.2in}{0.1mm}}\\[-1.3mm]
        {\committeenameC}        \hfill{Date}\\[2mm]
    \end{spacing}
  \end{minipage}
  \vfill
  
  {
	STEVENS INSTITUTE OF TECHNOLOGY\\
  	Castle Point on Hudson\\
  	Hoboken, NJ 07030\\
    {\graddate}
  }
  % \vfill

\end{center}
}
\makeatother

\newcommand{\thesiscopyrightpage}{
  \clearpage
  \thispagestyle{empty}      % no page number on the copyright page
    \strut \vfill
    \vspace*{7.2in}
   \begin{center}
        \textcopyright  2026, Daniel Kadyrov. All rights reserved.
    \end{center}
    \vfill \strut
  \clearpage}


\begin{document}

% \vspace*{-2\baselineskip}
\maketitle
\pagenumbering{roman}
\thispagestyle{empty}

\thesiscopyrightpage              

\chapter*{Abstract}
\addcontentsline{toc}{chapter}{Abstract}

Stored products, such as grains and processed foods, are a staple in global diets and one of the most traded commodities but are susceptible to infestation by various pests, most notably insects. The early detection of insects throughout the supply chain is crucial to mitigate against qualitative and quantitative losses, ensure food safety, and prevent introducing invasive species to new environments resulting in further economic and ecological impact.  

The Acoustic Stored Product Insect Detection System (A-SPIDS) was designed as a low-cost, non-destructive, alternative detection method that uses piezoelectric sensors to detect insect sounds and vibrations within an infested material. The system was used to record and analyze the vibro-acoustic signatures of various insects, including the \textit{Callosobruchus maculatus}, \textit{Tribolium confusum}, and \textit{Tenebrio molitor}, in different materials. These recordings were annotated and organized into a dataset accessible through a developed database software package.  

A normalization method was implemented using the ambient noise of the sensors as a reference to accommodate for proprietary, non-calibrated sensors and enable relative detection thresholds for unknown sensitives. The normalized envelope of the filtered signals was used to characterize and compare the insect signals by measuring the Normalized Signal Pulse Amplitude (NSPA) and Normalized Signal Energy Level (NSEL), parameters that apply the signal noise ratio (SNR) for pulse-based detection and averaged energy-based detection, respectively. 

The physical insulation in the platform was shown to provide an average attenuation of \SI{45}{\deci\bel} above \SI{2000}{\hertz}. The NSPA was applied as a detection metric resulting in a \SI{99.4}{\percent} detection accuracy at \SI{80}{\deci\bel} ambient noise levels when using a filter range of \SIrange{1565}{6000}{\hertz}. The external microphone integrated in the system design was leveraged to identify and remove insect-like signals from the piezoelectric sensor recordings. This led to a \SI{100}{\percent} detection rate with zero false alarms at noise levels exceeding \SI{100}{\deci\bel}.

\vfill

Author: Daniel Kadyrov\\
Advisor: Dr. Raju Datla and Dr. Alexander Sutin\\
Date: April 17, 2026\\
Department: Department of Civil, Environmental and Ocean Engineering\\
Degree: Doctor of Philosophy

\chapter*{Dedication}
\addcontentsline{toc}{chapter}{Dedication}

To my best friend and pup, Marx, and the thousands of insects that preceded him as pets and research subjects.

\chapter*{Acknowledgments}
\addcontentsline{toc}{chapter}{Acknowledgments}

First and foremost I would like to express my deepest gratitude to my advisor, Dr. Alexander Sutin, for his unwavering support, guidance, and encouragement throughout my Ph.D. journey. His expertise and wisdom in acoustics, signal processing, and the traversal of academic publications has been invaluable to my research, and his mentorship has helped me grow both academically and professionally.

I would like to thank my other committee members. Dr. Raju Datla for serving as the committee chair and guiding me throughout the Ph.D. process. Dr. Muhammad Hajj for his feedback that elevated the quality of my research. Dr. Richard Mankin for serving as an external member and the lengthy, extensive, and high quality research he has conducted in the field of insect detection that served as a foundation for this work. 

This journey would not have started if not for the initial motivation of my love, Gabriela Sanchez, and her enduring comfort despite the thousands of insects sharing our living space throughout the research. Marx was a puppy when I started the insect detection research and has been a steadfast and patient companion throughout the process, reminding me to take breaks and step outside. Thank you also to my family for their encouragement and my mother, Anya Braslavsky, for the technical drawings of the insect specimen. I am following in the footsteps of an aunt, grandfather, and step-father who all earned their Ph.D. degrees and have always wanted for me to do the same. 

Unquestionable support came from my former coworkers at the Sensor Technology and Applied Research (STAR) Center laboratory and Maritime Security Center (MSC) at Stevens Institute of Technology including Dr. Hady Salloum, Dr. Alexander Sedunov, Nikolay Sedunov, Christopher ``Kip'' Francis, Sergey Tsyuryupa, Aleksandr Merzhevskiy, Cameron Ruddy, Anthony Rusignioli, Beth DeFares, and Dr. Barry Bunin. Besides providing support with the development and deployment of systems and sensors, expanding my Russian lexicon to the dismay of my mother, and improving my chess game they all contributed to my growth as a researcher and engineering, instilling lifelong skills and passions as well as a strong importance and appreciation of engineering to national and global wellbeing.

The final appreciation goes to the Department of Homeland Security (DHS) Science and Technology Directorate (S\&T) for funding this research through the Cross-Border Threat Screening and Supply Chain Defense (CBTS) program under Grant Award Number 18STCBT00001. I was fortunate to join S\&T as a general engineer in the National Urban Security Technology Laboratory (NUSTL) and have received further support and encouragement from my colleagues there. 


% and providing the opportunity to work at the National Urban Security Technology Laboratory (NUSTL) during my Ph.D. studies, where I was able to continue this research and apply it to real-world problems. 
% The initial funding for this research was provided by the Cross-Border Threat Screening and Supply Chain Defense (CBTS) program through the Department of Homeland Security (DHS) Science and Technology Directorate (S\&T). Fortunately, I was able to join S\&T as a general engineer in the National Urban Security Technology Laboratory (NUSTL) during my Ph.D. studies, where. 


% Unquestionable support came from my former coworkers at the Sensor Technology and Applied Research (STAR) Center laboratory and Maritime Security Center (MSC) at Stevens Institute of Technology. Dr. Hady Salloum provided the intial opportunity to work at the STAR Center, where I was exposed to a variety of research topics and projects that guided my academic and professional journey. Dr. Alexander Sedunov and Nikolay Sedunov are two of the most brilliant minds I have ever met, and my research would not have been possible without their tireless hard work and expertise. My Russian lexicon was greatly expanded thanks to the colorful vocabulary of Sergey Tsur and Aleksandr Mershevsky. Christopher ''Kip`` Francis taught a curriculum never available in any university, how to correctly 


% Christopher ''Kip`` Francis

% The teachings and pedagogy of Dr. Barry Bunin extend far from the two courses I had the pleasure of spending with him. Along with Beth DeFares, Barry instilled the importance and appreciation of engineering to national and global wellbeing. 

% The teachings of Dr. Barry Bunin extend far from the two coursers 

% Dr. Barry Bunin 


% including Dr. Hady Salloum, Dr. Alexander Sedunov, Nikolay Sedunov, Christopher Francis, Sergey T., Aleksandr M., Cameron Ruddy, and Anthony Rusignioli. 



% \textbf{Committee Members}
% \begin{itemize}
% 	\item Dr. Alexander Sutin, Stevens Institute of Technology
% 	\item Dr. Muhammad Hajj, Stevens Institute of Technology
% 	\item Dr. Raju Datla, Stevens Institute of Technology
% 	\item Dr. Richard Mankin, USDA Agricultural Research Service
% \end{itemize}

% \newpage
% This research originated from a project funded by the Cross-Border Threat Screening and Supply Chain Defense (CBTS) to extend earlier work conducted by the Stevens Institute of Technology on the use of passive acoustic sensing for the detection of insects in stored products. The presentation and subsequent conference paper at the Acoustic Society of America Conference in 2022 on the topic of \citetitle{kadyrovImprovementsStevensDrone2022} showed promising results on the capabilities of the lab to accommodate detection of various targets using acoustics and signal-processing techniques. The preliminary work of research and development of alternative low-cost sensors to assist Custom and Border Protection Agricultural Specialists at United States Ports of Entry expanded into deeper analysis of utilizing the recordings made of live insect specimen in various materials using the sensor to study their acoustic characteristics, explore normalization methods, investigate reducing the influence of background noise, and eventually implement a detection algorithm.

% In \citetitle{kadyrovVibroAcousticSignaturesVarious2024}, we introduced the Acoustic Stored Product Insect Detection System (A-SPIDS) and published a comprehensive analysis of the signatures of various stored product insects in different materials based on the recordings. Highlighting their unique spectral and temporal characteristics, we developed normalization methods that used the sensor's self-noise to standardize the recordings between different, non-calibrated, piezoelectric discs deployed in the sensor. This normalization method was extended to use the median of the signal and shown that it could be used when a sensor's self-noise data is not available, like in the case of the publicly available BugBytes dataset. 

% These findings were presented as a poster ``Perspectives of Machine Learning for Acoustic Methods of Stored Product Insect Detection''  at the Artificial Intelligence Meets Science (AIMS) Conference in 2024. After publishing the underlying data and software from the recordings as the Acoustic Stored Product Insect Dataset and Acoustic Stored Product Insect Database on Kaggle and GitHub, respectively, we received interest from the NYU AI for Scientific Research (AIfSR) group to explore applying machine learning techniques on the data. Although there were interesting findings, we were not able to publish in NeurIPS because the recording subjects were not sufficiently diverse. A presentation was also made at the DHS RDT\&E summit.

% In the upcoming publication, titled ``Acoustic Detection of Insects in Stored Products in the Presence of Strong Ambient Noise", we will focus on how we could use the spectral and temporal characteristics, normalization, and metrics we found in the earlier publication to contrive a robust detection algorithm. The article begins by showing that although the insulation of the A-SPIDS sensor provides some protection against ambient noise, it is not entirely effective and often appears as similar signals to those produced by insects. At this moment, we have made significant progress in showing how the normalization metrics could be used in order to optimize the examination frequency range to avoid this interference. The next steps are to implement the detection algorithm and evaluate its performance.

% table of contents 
\newpage
\setcounter{tocdepth}{1}
\begingroup
\setlength{\parskip}{0pt}
\setstretch{1} % distance between the sections
\small
\tableofcontents
\endgroup
\newpage
\begingroup
\setlength{\parskip}{0pt}
\setstretch{1}
\small
\listoftables
\endgroup
\newpage
\begingroup
\setlength{\parskip}{0pt}
\setstretch{1}
\small
\listoffigures
\endgroup


\newpage
\pagenumbering{arabic}
% \onehalfspacing
% \setstretch{1.125}


\subfile{1_introduction/introduction}
\subfile{2_literaturereview/literature_review}
\subfile{3_methodology/methodology}
\subfile{4_characteristics/characteristics}
% \subfile{5_results/results}
\subfile{5_results/detection}
\subfile{6_conclusion/conclusion}

\newpage
{\small\singlespacing
	\printbibliography[heading=bibintoc]
}

\newpage
\appendix
\appendixpage
\subfile{appendix/a1_noise_spectrograms}
\subfile{appendix/a2_scripts}

\newpage
% Remove chapter numbering for CV
\chapter*{Curriculum Vitae}
\addcontentsline{toc}{chapter}{Curriculum Vitae}
% Force single spacing and remove all extra spacing
\begin{singlespace}
\footnotesize % Make font smaller for CV
\setlength{\parindent}{0pt} % no indent for CV
\setlength{\parskip}{0pt}   % no space between paragraphs
\setlength{\itemsep}{0pt}   % no space between items
\setlength{\parsep}{0pt}    % no space between paragraphs in lists
\setlength{\topsep}{0pt}    % no space above/below lists
\setlength{\partopsep}{0pt} % no extra space above/below lists

\subsection*{Daniel Kadyrov}

Email: \href{mailto:dkadyrov@stevens.edu}{dkadyrov@stevens.edu}\\

\subsection*{Education}
\textbf{Stevens Institute of Technology} \hfill Hoboken, NJ \\
Ph.D. Ocean Engineering \hfill Expected Spring 2026 \\
M.S. Computer Science (Machine Learning and Data Science) \hfill May 2021 \\

\textbf{Binghamton University, State University of New York} \hfill Binghamton, NY \\
M.S. Mechanical Engineering (Acoustics and Dynamic Systems) \hfill May 2016 \\
B.S. Mechanical Engineering \hfill May 2015 \\

\subsection*{Professional Experience}

\textbf{Department of Homeland Security}\\
National Urban Security Technology Laboratory (NUSTL) \hfill New York, NY \\
\textit{General Engineer \hfill May 2019 - Present} \\

\textbf{Stevens Institute of Technology}\\
Sensor Technology and Applied Research Center \hfill Hoboken, NJ \\
\textit{Research Engineer \hfill May 2019 - December 2023} \\

\textbf{Thornton Tomasetti}\\
Weidlinger Applied Science \hfill New York, NY \\
\textit{Senior Engineer \hfill January 2017 - November 2018} \\

\section*{Teaching Experience}
\textbf{Pace University} \hfill New York, NY \\
\textit{Adjunct Professor \hfill September 2023 - Present} \\

\textbf{Columbia University} \hfill New York, NY \\
\textit{Adjunct Lecturer \hfill June 2023 - August 2024} \\

\subsection*{Journal Publications}
\begin{itemize}
  \item D. Kadyrov, A. Sutin, N. Sedunov, A. Sedunov, and H. Salloum, “Acoustic Detection of Insects in Stored Products in the Presence of Strong Ambient Noise,” Sensors, Feb. 2026, doi: 10.3390/s26051511.
  \item D. Kadyrov, A. Sutin, N. Sedunov, A. Sedunov, and H. Salloum, “Vibro-Acoustic Signatures of Various Insects in Stored Products,” Sensors, vol. 24, no. 20, Oct. 2024, doi: 10.3390/s24206736.
\end{itemize}

\subsection*{Conferences}
\begin{itemize}
  \item DHS Research Development \& Testing and Evaluation Summit (Online, 2025)\\
  \textit{Stored Product Insect Detection Systems (SPIDS)}
  \item Artificial Intelligence Meets Science (AIMS) Conference (New York, 2024)\\
  \textit{Perspectives of Machine Learning for Acoustic Methods of Stored Product Insect Detection}
  \item Acoustic Society of America (Denver, 2022)\\
  \textit{D. Kadyrov, A. Sedunov, N. Sedunov, A. Sutin, H. Salloum, and S. Tsyuryupa, “Improvements to the Stevens drone acoustic detection system,” Proc. Mtgs. Acoust., vol. 46, no. 1, p. 045001, Aug. 2022, doi: 10.1121/2.0001602.}
\end{itemize}
\end{singlespace}



\end{document}
